\subsection*{MT-M14. Bounded multiplicities and the full Hermite basis}
Retain the canonical completion and quotient of MT-M12, assume
GRH, and replace simplicity by the hypothesis
\[
                  1\le m_\rho\le m_0<\infty
                              \qquad(\rho\in Z),       \tag{13}
\]
where $m_\rho$ is the actual zero multiplicity.
Then the following are equivalent:
\begin{enumerate}
\item For some $c,C>0$,
$|L^{(m_\rho)}(\rho,\eta)/m_\rho!|
       \ge c(1+|\operatorname{Im}\rho|)^{-C}$ for all $\rho$.
\item The canonical quotient has an absolute Schauder basis
whose coefficient functionals are
$[f]\mapsto(\mathcal Mf)^{(k)}(\rho)/k!$ for all
$0\le k<m_\rho$.
\item It has an ordinary Schauder basis with precisely those
continuous coefficient functionals.
\end{enumerate}
Under these conditions the closed ideal is the full zero-jet
kernel, and the quotient is the rapid sequence space on
$\{(\rho,k):0\le k<m_\rho\}$, weighted by
$1+|\operatorname{Im}\rho|$. Assumption (13) and condition (1)
are additional arithmetic hypotheses, not consequences of GRH
established in this run.

\begin{proof}
The necessity of (1) from (3) is exactly MT-M11.
We give all changes needed for sufficiency; they also show why
the fixed bound $m_0$ is used.
Write $z=s-\rho$, $m=m_\rho$, and
\[
                  G_\rho(s)=\Xi_\eta(s)/(s-\rho)^m .
\]
The leading-coefficient hypothesis, multiplied by the gamma
and conductor factors, gives
\[
 |G_\rho(\rho)|\ge c(1+|T|)^{-C}e^{-\kappa|T|},
           \qquad T=\operatorname{Im}\rho.
\]
Cauchy's formula on a fixed larger disk around $\rho$ and
the strip bound (5) give
$|G_\rho^{(j)}(\rho)/j!|
 \le C(1+|T|)^C e^{-\kappa|T|}$ for $0\le j\le m_0$.
To see the division estimate directly, on a circle $|z|=2$
divide $\Xi_\eta(\rho+z)$ by $z^m$ and apply Cauchy's
formula inside; $m\le m_0$ keeps all constants fixed.

Fix $a>0$. Let $P_{\rho,k}(z)$ be the Taylor polynomial
through degree $m-1-k$ of the analytic germ
$e^{-az^2}/G_\rho(\rho+z)$.
The reciprocal-coefficient recursion starts with
$1/G_\rho(\rho)$ and at each subsequent degree divides
a finite sum by $G_\rho(\rho)$. The displayed upper and
lower bounds therefore bound every coefficient of
$P_{\rho,k}$ by
$C(1+|T|)^C e^{\kappa|T|}$, uniformly for
$0\le k<m\le m_0$. Only finitely many recursion steps
are involved. Define
\[
 H_{\rho,k}(s)=
      z^kP_{\rho,k}(z)e^{az^2}G_\rho(s).                \tag{14}
\]
Its Taylor expansion at $\rho$ is
$z^k+O(z^m)$, and it vanishes to the required multiplicity
at every other zero. On every fixed vertical strip its
bound has the same Gaussian form (8), with constants
depending also on $m_0$: away from $|z|<1$, use (5)
and the fixed polynomial degree in (14); inside,
use the Cauchy bound for $G_\rho$.
Inverse Mellin transforms $h_{\rho,k}\in\mathcal E$
therefore satisfy polynomial seminorm bounds (9)
uniformly over $k<m_\rho$.

The coefficients $(\mathcal Mf)^{(k)}(\rho)/k!$ are
uniformly rapidly decreasing in $|T|$ for
$k<m_0$. Indeed, they are Mellin transforms of the
fixed finite list $x^k f(x)/k!$, and the
integration-by-parts argument of MT-M12 applies to
each. The series with coefficients these jets and
vectors $h_{\rho,k}$ consequently converges absolutely
in $\mathcal E$, is continuous in $f$, and interpolates
all its jets. Conversely every rapid sequence of jets
is obtained from the corresponding convergent series.
Multiplicity is included in the local zero count,
so the same summability estimate applies.

It remains to verify the full ideal equality.
Repeat the disk argument proving (11), now with
Blaschke factors raised to the actual multiplicities.
The number of factors counted with multiplicity is still
$M=O_A(\log(2+|t|))$. If a nearest zero $\rho_0$
is in the inner half disk, the leading coefficient
identity is
\[
 \frac{L^{(m_{\rho_0})}(\rho_0)}{m_{\rho_0}!}
   =G(\rho_0)
       \lim_{s\to\rho_0}
                    \frac{B(s)}{(s-\rho_0)^{m_{\rho_0}}}.
\]
The upper bound for that leading Blaschke coefficient
is $C_A^M$ times
$\prod_{\rho_j\ne\rho_0}
          |\rho_0-\rho_j|^{m_{\rho_j}}$.
Condition (1) therefore supplies the polynomial lower
bound (12) for this multiplicity-weighted product.
Nearest-zero comparison gives
\[
 |L(s,\eta)|\ge c_A(1+|t|)^{-C_A}
                         |s-\rho_0|^{m_{\rho_0}}
\]
when the nearest distance is less than one.
When that distance is at least one, or the nearest
disk zero is outside the inner half disk, the same
argument gives a polynomial lower bound without
the distance factor. All constants are uniform
because $m_{\rho_0}\le m_0$ and $M=O_A(\log|t|)$.

If $F=\mathcal Mf$ vanishes to the required orders,
Taylor's formula at a nearest zero gives
\[
 F(s)=\frac{(s-\rho)^m}{(m-1)!}
       \int_0^1(1-u)^{m-1}
                         F^{(m)}(\rho+u(s-\rho))\,du .
\]
For $m\le m_0$ its derivative bound is uniformly
rapid in $|t|$. The distance factors cancel.
Thus $F/L$ is rapidly decreasing on each fixed
strip at large height. Exactly the Gaussian-regularized
division and approximate-identity proof in MT-M12
shows that $f$ lies in
$\overline{\psi_\eta*\mathcal E}$.
The reverse inclusion is immediate.
The interpolation series now identifies the quotient
topologically with the rapid jet-sequence space,
as in MT-M12. This proves (2), hence (3), and
completes the equivalence.
\end{proof}

No unspecified infinite Hermite estimate has been hidden in
the argument: the reciprocal-coefficient recursion has at most
$m_0$ steps. With unbounded multiplicities, the exponents of
the polynomial bounds produced by that recursion need not be
uniform. That case remains a separate interpolation gate.

\subsection*{MT-M15. Multiple zeros: positive traces and nilpotent operators}
Under the conditions of MT-M14, convolution has an explicit
nuclear representation on the rapid jet-sequence quotient
and a trace-class realization on the standard $\ell^2$
space of those jets, with
\[
          \operatorname{Tr}C_f
                    =\sum_{\rho\in Z}m_\rho\mathcal Mf(\rho).
                                                               \tag{MT-M15}
\]
For $f=h*h^\sharp$ this trace is
$\sum_\rho m_\rho|\mathcal Mh(\rho)|^2\ge0$.
However, if any $m_\rho>1$, there is no positive-definite
inner product on the full faithful jet quotient for which
every such adjoint-square convolution is a positive
operator. In particular, the trace positivity does not
upgrade to positivity of the entire nonsemisimple action.

\begin{proof}
At a zero $\rho$ let $a_j(f)=(\mathcal Mf)^{(j)}(\rho)/j!$.
The product rule in normalized jets is
\[
       a_k(f*g)=\sum_{j=0}^k a_j(f)a_{k-j}(g).
\]
Thus the block of $C_f$ at $\rho$ is a lower triangular
Toeplitz matrix with diagonal $a_0(f)$ and size $m_\rho$.
All its entries decrease rapidly with the ordinate,
uniformly in the at most $m_0$ indices.
The trace norm of a finite matrix is at most the sum
of the trace norms of its individual matrix units,
so the block trace norms are summable. This proves
the Hilbert trace-class assertion and its trace.

For nuclearity, use the rank-one expansion in these
matrix entries. The jet functionals are equicontinuous:
for a fixed $N>3/2$,
\[
 |(\mathcal Mg)^{(k)}(\rho)/k!|
 \le \frac1{k!}\int_{\mathbb R}|x|^k e^{x/2}|g(x)|\,dx
 \le2p_{N,0}(g),
\]
uniformly even before using $k<m_0$.
The interpolation vectors have polynomial seminorm
bounds, and the matrix entries have rapid decay.
Insert the summable scalar weight $(1+|T|)^{-2}$
as in MT-M13. With at most $m_0^2$ matrix units per
zero, this is an explicit nuclear representation.
The weighted involution has jet coefficients
\[
                   a_j(h^\sharp)=(-1)^j\overline{a_j(h)}.
\]
Its zeroth coefficient gives the asserted nonnegative
trace for an adjoint square.

Finally suppose $m=m_\rho>1$. The interpolating functions
in MT-M14 realize every local jet and zero jets elsewhere.
The local algebra is $\mathbb C[z]/(z^m)$, with
$z^\sharp=-z$ and conjugation on scalars.
Let $p$ denote its identity, extended by zero in every
other zero block, and let $n=z$ in this block and zero
elsewhere. Both are represented by actual functions
of $\mathcal E$, by (14). On the $m$-dimensional local
block, their convolution operators are $I$ and a
nonzero nilpotent $N$.
For every real $u$, the class $h=p+iu\,n$ satisfies
$h^\sharp=h$, and hence its adjoint square acts there as
\[
                         I+2iuN-u^2N^2.
\]
If all adjoint-square operators were positive in a
positive-definite inner product, these operators and
the corresponding operators with $-u$ would be
self-adjoint on the invariant local block.
Subtracting them for $u\ne0$ shows that $iN$ is
self-adjoint, equivalently $N^*=-N$.
A skew-adjoint operator on a finite-dimensional positive
inner-product space is normal and diagonalizable;
nilpotence forces all eigenvalues to vanish, so it is
zero. This contradicts the regular action of $z$
on $\mathbb C[z]/(z^m)$.
The argument is internal to a finite block and
does not depend on any choice of Hilbert completion.
\end{proof}

Passing to the semisimple quotient that retains only the
zeroth coordinate of each block removes the nilpotents.
Using $m_\rho$ copies of each scalar coordinate then
recovers the positive spectral trace with multiplicity.
It does not faithfully recover the discarded higher
jets. This is the precise distinction needed when
transferring Alpöge-style positivity to an orbital
jet quotient in Masullo 9.3.2.5--9.3.3, or to the
cote~1/SA23 trace programme. No assertion that the
draft explicitly claims a positive faithful Hilbert
realization of this entire quotient is made here.
